\begin{abstract}
The generalized lambda distribution (GLD) is a four-parameter quantile-defined
family able to reproduce a wide range of shapes with a single functional form,
which makes it attractive whenever a simulation output must be summarised by a
distribution rather than by a few moments. Its software support, however, is
concentrated in R, and the Python packages available expose one parameterization
at a time and no facility for assessing how good a fit actually is. This paper
presents \texttt{pyGLAM}, an open-source Python package implementing the FKML and
the Ramberg-Schmeiser parameterizations behind a single interface, fitting them
by multi-start moment matching restricted to valid distributions, and scoring
the result with a goodness-of-fit module. Its quantile, distribution and density
functions and its moments agree with the R packages \texttt{gld} and
\texttt{GLDEX} to numerical precision, and its fits of the same samples are
valid, reproduce the sample moments and run several times faster. A factorial
Monte Carlo study over two parameterizations, six target distributions and four
sample sizes, with one hundred replicates per cell, measures the error of each
fit against the analytical target. The slope of the Kolmogorov-Smirnov distance
against sample size, with a bootstrap confidence interval, separates a sampling
component, which decays at the rate $N^{-1/2}$, from a structural component,
which does not. The normal and the uniform show no structural component, the
Gumbel and the Beta a small one, and the Gamma a floor under both
parameterizations, which match the same moments with distributions that fail in
different places. On a Student-$t$ with four degrees of freedom the same
diagnostic points to the estimator rather than the family. The study also shows
that moment matching admits roots that are not distributions and are reported as
converged by the optimizer, and that the package keeps every selected fit valid.
\end{abstract}

\noindent\textbf{Keywords}: generalized lambda distribution; quantile
distributions; method of moments; goodness of fit; statistical software; Python.
